\section*{Complete author Section 5.1 proof and required mathematical tools}
A formal power series is a generalization of a polynomial, where the number of terms can be infinite. Formally, $A=\sum_{i\ge0}A_ix^i$ with $A_i\in\mathbb F$, is a power series in the power series ring $\mathbb F[[x]]$. We define the degree $d$ truncation $\mathrm{trunc}(A,d)$ of $A$ to be $\mathrm{trunc}(A,d):=\sum_{0\le i\le d}A_ix^i$. So, $\mathrm{trunc}(A,d)$ is always a polynomial of degree at most $d$.
We show that complexity of truncation of rational functions where the degrees are small, has low complexity. For simplicity, we work with $\mathbb F=\overline{\mathbb F}$, an algebraically closed field. We first recall the following folklore decomposition.
\paragraph{Lemma 24 (Partial fraction decomposition).} Let $g(x)/h(x)$ be a rational function with $\deg(g)<\deg(h)$. If $h(x)=\prod_{i\in[k]}(x-a_i)^{d_i}$ is the factorization of $h(x)$ over $\mathbb F[x]$, then, there exist $b_{ij}\in\mathbb F$ s.t.:
\[g(x)/h(x)=\sum_{i\in[k]}\sum_{j\in[d_i]}b_{ij}/(x-a_i)^j.\]
\paragraph{Lemma 25.} For any non zero $a\in\mathbb F$, we have $L(\mathrm{trunc}(1/(x-a),d))=O(\log d)$.
\begin{proof}This follows from the inverse identity: $1/(a-x)=1/a\sum_{i\ge0}(x/a)^i$ and the fact that $L(\sum_{0\le i\le d}(x/a)^i)=O(\log d)$ (By using Remark 15).\end{proof}
\paragraph{Lemma 65 (Power series valuation).} Let $g,h\in\mathbb F[x]$ such that $g/h\in\mathbb F[[x]]$. Let $m$ (respec. $n$) be the highest power dividing $g$ (respec. $h$) i.e. $x^m\mid g$ and $x^{m+1}\nmid g$ (respec. for $h$). Then, $m\ge n$.
\begin{proof}Suppose, $m<n$. Note that, there exists $0\ne\alpha\in\mathbb F$, such that $h=\alpha x^n(1+x\widetilde h)$, for some $\widetilde h\in\mathbb F[x]$. Similarly, let $g=\beta x^m(1+x\widetilde g)$, for some $\widetilde g\in\mathbb F[x]$ and $\beta\in\mathbb F$. Thus,
\begin{align*}\frac gh&=\frac\beta\alpha x^{m-n}\frac{1+x\widetilde g}{1+x\widetilde h}\\&=\frac\beta\alpha x^{m-n}(1+x\widetilde g)(1+x\widetilde h+(x\widetilde h)^2+\cdots)\\&\notin\mathbb F[[x]],\text{ a contradiction}.\end{align*}\end{proof}
\paragraph{Lemma 66 (General Leibniz rule).} If $f$ and $g$ are $k$-time differentiable functions, then $(fg)^{(k)}=\sum_{i=0}^k\binom ki f^{(k-i)}g^{(i)}$.
\paragraph{Lemma 63.} Let $f$ be a polynomial $\mathbb F[\mathbf x]$, computed by a size $s$ circuit $C$. Then, there exists a circuit $C'$ of size $O(sm^2)$ which computes $f,f^{(1)},f^{(2)},\ldots,f^{(m)}$.
\begin{proof}We split every $G$ gate in $C$ to $n+1$ gates $G_0,\ldots,G_m$ in $C'$. The property we want is that if the gate $G$ is computing the polynomial $g$ in $C$ then $G_k$ computes the polynomial $g^{(k)}$ in $C'$. Suppose $G$ is a $+$ gate in $C$ with children gates computing the polynomials $g_1$ and $g_2$. Now we know that $g^{(k)}=g_1^{(k)}+g_2^{(k)}$. Thus we can easily propagate the derivatives on addition/subtraction gates. If $G$ is a $\times$ gate then using Lemma 66, we know that:
\[(g_1g_2)^{(k)}=\sum_{i=0}^k\binom ki g_1^{(k-i)}g_2^{(i)}.\]
Thus we can computes $g,g^{(1)},g^{(2)},\ldots,g^{(m)}$ using additional $O(m^2)$ gates. Therefore $C'$ has $O(sm^2)$ gates.\end{proof}
\paragraph{Proof of Theorem 26.} The main idea is to use Lemma 24 and the low complexity of the truncation of inverse identity (Lemma 25). However, the given polynomial $h$ may be divisible by $x$ (i.e. $h(0)=0$). In that case, let $m$ be the highest power of $x$ such that $x^m\mid h$ (i.e. $x^{m+1}\nmid h$). Note that, as $g/h\in\mathbb F[[x]]$, $x^m\mid g$ as well (Lemma 65). As $\deg(h)\le d_h$, thus $m\le d_h$.
By using Theorem 18, we know that $g_1:=g/x^m$ has a cicuit of size $O((s+\log d_h)d_h^2)=:s_1$. Trivially, $h_1:=h/x^m$ has degree $\le d_h$, and $g/h=g_1/h_1$. Denote, $g_2:=g_1\mathbin{\mathrm{mod}}h_1$. Obviously, $\deg(g_2)<\deg(h_1)$ and $g_1/h_1=g_1\mathbin{\mathrm{div}}h_1+g_2/h_1$. Invoking Theorem 13, one concludes that $L(g_1\mathbin{\mathrm{div}}h_1)=O(s_1d_h)$. Therefore, $L(\mathrm{trunc}(g_1\mathbin{\mathrm{div}}h_1,d))=O(s_1d_h)$, as $\deg(g_1)<d$.
Let $h_1$ factors over $\mathbb F[x]$ as $h_1:=\prod_{i\in[k]}(x-a_i)^{d_i}$. Trivially, $\sum d_i\le d_h$. By using Lemma 24 on $g_2/h_1$, we know that there are constants $a_i,b_{ij}\in\mathbb F$ such that:
\[g_2(x)/h_1(x)=\sum_{i\in[k]}\sum_{j\in[d_i]}b_{ij}/(x-a_i)^j.\]
Note that, for any $a\in\mathbb F$ and $t\in\mathbb N$, $d^t/dx^t(1/(x-a))=(-1)^tt!\cdot1/(x-a)^{t+1}$, and thus,
\begin{multline*}\mathrm{trunc}(1/(x-a)^{t+1},d)\\=\frac{(-1)^t}{t!}\frac{d^t}{dx^t}(\mathrm{trunc}(1/(x-a),d))+\sum_{i=d-t+1}^d\gamma_ix^i,\qquad\gamma_i\in\mathbb F.\end{multline*}
Using the above identity and Lemma 63, we can show that
\[L\left(\mathrm{trunc}\left(\sum_{j\in[d_i]}b_{ij}/(x-a_i)^j,d\right)\right)=O(\log d\cdot d_i^2).\]
To show this, note that $L(\mathrm{trunc}(1/(x-a_i),d))=O(\log d)$, and using Lemma 63, we compute all its derivative till the $d_i$-th one which has a circuit of size $O(\log d\cdot d_i^2)$. Using the above identity, we can add $d_i-1$ many monomials of the form $cx^\ell$ with $d-d_i+2\le\ell\le d$ (each monomial has trivial size of $O(\log d)$) to the circuit to obtain a circuit for $\mathrm{trunc}(\sum_{j\in[d_i]}b_{ij}/(x-a_i)^j,d)$, which still has size $O(\log d\cdot d_i^2)$. Thus, doing it for each $a_i$ for $i\in[k]$, one obtains that
\begin{align*}L(\mathrm{trunc}(g(x)/h(x),d))&=L(\mathrm{trunc}(g_1(x)/h_1(x),d))\\&=L(g_1\mathbin{\mathrm{div}}h_1,d)+L(g_2/h_1,d)\\&=O(s_1d_h)+L\left(\mathrm{trunc}\left(\sum_{i\in[k]}\sum_{j\in[d_i]}b_{ij}/(x-a_i)^j,d\right)\right)\\&=O((s+\log d_h)d_h^3)+O\left(\log d\cdot\sum_{i\in[k]}d_i^2\right)\\&=O(sd_h^3\log d).\end{align*}\qed
\paragraph{Remark 27.} Eventually, we can replace $g\in\mathbb F[[x]]$ with the given complexity $\mathrm{trunc}(g,d)=s$ and show that the exact same proof as above, works.
